\section{Introduction}

High-dimensional geometric structures have become increasingly relevant in the design of multimodal cognitive–narrative control systems. Such systems require mathematical tools capable of expressing both global manifold behavior and local curvature-sensitive dynamics. This preprint focuses on three foundational components that enable this interaction: (1) 80-dimensional manifold embeddings, (2) 9-dimensional radial Laplacian operators, and (3) Σ/Λ save-state geometry.

The 80D embedding provides a global geometric environment in which world manifolds, cognitive states, and narrative structures can be represented with stability envelopes and holonomy-regulated transport. The 9D radial Laplacian serves as a local differential operator acting on these embedded manifolds, capturing curvature response, modal adjacency, and spectral diffusion. Σ/Λ save-state geometry introduces spectral and holonomy snapshots that allow drift-neutral restoration of cognitive and narrative configurations.

Although these components arise naturally within broader systems such as NDH cognitive geometry, FSFE math kernel, and Soft-World manifold architectures, the present work isolates and formalizes the mathematical core. By focusing on the embedding map, the radial Laplacian, and the save-state operators, we establish a unified operator–manifold coupling that can serve as a foundation for future geometric AI research.

The remainder of this preprint is organized as follows. Section 2 introduces the mathematical preliminaries, including manifolds, metrics, tensor fields, and differential operators. Section 3 formalizes the 80D embedding and its induced geometric structures. Section 4 derives the 9D radial Laplacian and analyzes its curvature-sensitive behavior. Section 5 introduces Σ/Λ save-state geometry and its stability conditions. Section 6 presents the operator–manifold coupling between the Laplacian and the embedded manifold. Sections 7 and 8 provide discussion and conclusions. Appendices contain formal definitions, theorems, proofs, commutation relations, and diagrams.
